\section{Mass Model and Likelihood}
\label{sec:model}
\label{app:model}
\label{sec:bayes_model}
\emph{In short: a Gaussian signal and a double-exponential background for each charge, with the raw asymmetry $A$ as a parameter of the fit; both inference methods use the same likelihood.}
The mass spectrum of each charge sample is described by a signal-plus-background mixture on the nominal fit window $[5100,5600]\,\MeV$ (Sec.~\ref{sec:data_mass}),
\begin{equation}
f^{\pm}(m)=N_s^{\pm}\,S(m;\mu,\sigma)+N_b^{\pm}\,B^{\pm}(m),
\label{eq:mixture}
\end{equation}
where $S$ is a Gaussian truncated to the window, with mean $\mu$ and width $\sigma$ shared by the two charges, and the background is a double exponential with a shape of its own for each charge,
\[
B^{\pm}(m)=f_{\pm}\,E(m;\lambda_{1}^{\pm})+(1-f_{\pm})\,E(m;\lambda_{2}^{\pm}),
\]
with $E(m;\lambda)\propto e^{-\lambda m}$ a truncated exponential normalized on the window. The steep component ($\lambda_1$) describes the rise of the background towards the lower edge of the window, the slow one ($\lambda_2$) the combinatorial background under and above the peak. The ordering $\lambda_1>\lambda_2$ removes the label-switching symmetry of the mixture. Signal and background yields are split between the charges as
\[
N_s^{\pm}=\tfrac{N_s}{2}(1\mp A),\qquad N_b^{\pm}=\tfrac{N_b}{2}(1\mp A_b),
\]
so that $A=(N_s^--N_s^+)/(N_s^-+N_s^+)$, the raw asymmetry \Araw of Eq.~\eqref{eq:araw}, is itself a parameter of the fit, the parameter of interest (POI, all the others being nuisance parameters), and the background asymmetry $A_b$ is free. The parameter vector has twelve components, $\theta=(N_s,N_b,A,A_b,\mu,\sigma,\lambda_1^{\pm},\lambda_2^{\pm},f_{\pm})$.
Given the data $D=\{(m_i,q_i)\}_{i=1}^{N}$ (mass and charge of each candidate), the extended unbinned likelihood of the two charge samples is
\begin{multline}
\mathcal{L}(D\mid\theta)\propto e^{-(N_s+N_b)}
\prod_{i:\,q_i=+1}\Big[N_s^{+}S(m_i)+N_b^{+}B^{+}(m_i)\Big] \\
\times\prod_{i:\,q_i=-1}\Big[N_s^{-}S(m_i)+N_b^{-}B^{-}(m_i)\Big],
\label{eq:likelihood}
\end{multline}
where $S(m_i)\equiv S(m_i;\mu,\sigma)$. The likelihood is extended (the yields are parameters, hence the Poisson term, Sec.~\ref{app:model}); it is the common ground of the two mass analyses: Sec.~\ref{sec:frequentist} maximizes it and Sec.~\ref{sec:bayesian} multiplies it by a prior. Since both use this model unchanged, any difference between their results can only come from the inference.
\subsection{Extended Likelihood}
\emph{In short: with a Poisson total number of candidates the likelihood is $e^{-(N_s+N_b)}\prod_i[N_s S+N_b B]$ up to a constant, which constrains the yields and $A$ together.}
If the total number of candidates $N$ is Poisson-distributed with mean $\nu=N_s+N_b$ and each candidate is drawn independently from the normalized density $f(m,q)=[N_s S+N_b B]/\nu$, the likelihood of the sample is $\frac{\nu^N e^{-\nu}}{N!}\prod_{i}f(m_i,q_i)=\frac{e^{-\nu}}{N!}\prod_i[N_s S(m_i)+N_b B(m_i)]$. The factor $1/N!$ does not depend on the parameters and is dropped, which gives the extended likelihood of Eq.~\eqref{eq:likelihood}. Unlike the likelihood with a fixed total, it constrains the yields $N_s$ and $N_b$ through $N$ itself, and the asymmetry $A$ is estimated together with them.
\subsection{Extended Likelihood versus a Mixture with Fixed Total}
\emph{In short: the extended form differs from a mixture with fixed total only by the Poisson term of the total yield, which carries no information on $A$; the purity of the signal window is about $96\%$.}
The extended form differs from a mixture of signal and background with a fixed total number of candidates and a signal fraction only through the Poisson term for the total yield, which carries no information on $A$ and does not change its estimate; it is used because the yields of the two charges are the natural parameters of the model. Bayes' theorem applied to the fitted yields gives the probability that a candidate in the signal window $|m-\mu|<2\sigma$ is a signal decay, about $96\%$: a purity, not a property of the single candidate, which depends on the signal fraction of the sample (here $78\%$ of the candidates in the fit window).
\subsection{Label Switching of the Background Components}
\emph{In short: the mixture is invariant under the exchange of the two exponentials, which would give two identical modes; the ordering $\lambda_1>\lambda_2$ removes the symmetry.}
The mixture is invariant under the exchange $(\lambda_1,f)\leftrightarrow(\lambda_2,1-f)$, which would give the posterior two identical modes (label switching); the ordering $\lambda_1>\lambda_2$ removes this symmetry.
\subsection{Minimal Model as a Cross-Check}
\label{sec:bayes_minimal}
\emph{In short: a six-parameter model on a narrower window gives $\Araw=-0.0497$ (68\% CI $[-0.058,-0.041]$), compatible with the nominal result within its uncertainty; it is the model used for the low-statistics comparison.}
Before the nominal model was fixed, the same analysis was carried out with a minimal six-parameter model: a narrower window, $[5150,5450]\,\MeV$, which excludes most of the rising background below the peak, a single exponential background with a shape and a yield shared by the two charges, and the same Gaussian signal. On its $17{,}839$ candidates it gives $\Araw=-0.0497$ (68\% CI $[-0.058,-0.041]$) with $\chi^2/\text{ndof}=127.4/54$, and $-0.0457$ when the background yields of the two charges are left free. Its toy validation, with 500 toy-MC datasets, gives coverage consistent with nominal (0.664 and 0.962 for the 68\% and 95\% intervals). The difference from the nominal result, $0.005$, is driven by the background treatment, as in the systematic study of Sec.~\ref{sec:syst}, and is smaller than the statistical uncertainty.
\FloatBarrier
